\clearpage
\section{原题索引}
本索引按原稿题号排列。页码指本书页脚编号，可直接点击跳转；正文把题目放在所需知识之后，以便从概念连续学到应用。
\begin{longtable}{p{13mm}p{119mm}p{17mm}}
\toprule 题号 & 主题 & 页码\\\midrule
\endhead
A1 & 四个掩模的逐项判读 & \pageref{original:A1}\\
A2 & 三个滤波核的性质和方向 & \pageref{original:A2}\\
B1 & 四类直方图的明暗与对比度 & \pageref{original:B1}\\
B2 & 夜间道路双峰直方图 & \pageref{original:B2}\\
B3 & 离散均衡为何不完全均匀 & \pageref{original:B3}\\
B4 & 频谱图判读 & \pageref{original:B4}\\
C1 & 暗人物与亮屏幕的直方图 & \pageref{original:C1}\\
C2 & 灰度变换曲线 & \pageref{original:C2}\\
D1 & 中值滤波五个问题 & \pageref{original:D1}\\
D2 & 均值与中值的选择 & \pageref{original:D2}\\
D3 & 滤波核系数和 & \pageref{original:D3}\\
D4 & 拉普拉斯与一、二阶差分 & \pageref{original:D4}\\
D5 & 非锐化掩模与高提升 & \pageref{original:D5}\\
D6 & 完整频域滤波过程 & \pageref{original:D6}\\
D7 & 同态滤波 & \pageref{original:D7}\\
E1 & 直方图均衡完整计算 & \pageref{original:E1}\\
E2 & 相机视场与物方分辨率 & \pageref{original:E2}\\
E3 & 平均去噪的方差与帧数 & \pageref{original:E3}\\
E4 & 韦伯比 & \pageref{original:E4}\\
E5 & 采样周期与采样间隔 & \pageref{original:E5}\\
E6 & 频域滤波数值及振铃 & \pageref{original:E6}\\
F1 & 增强与复原 & \pageref{original:F1}\\
F2 & 低级与高级处理 & \pageref{original:F2}\\
F3 & 视锥与视杆 & \pageref{original:F3}\\
F4 & 采样与量化 & \pageref{original:F4}\\
F5 & 空间与灰度分辨率 & \pageref{original:F5}\\
F6 & 邻域与混合邻接 & \pageref{original:F6}\\
F7 & 算术与逻辑运算 & \pageref{original:F7}\\
G1 & 报靶系统设计 & \pageref{original:G1}\\
G2 & 隧道裂缝检测 & \pageref{original:G2}\\
G3 & 车牌增强方案 & \pageref{original:G3}\\
H1 & 椒盐噪声与中值 & \pageref{original:H1}\\
H2 & 点处理 & \pageref{original:H2}\\
H3 & 二阶差分双响应 & \pageref{original:H3}\\
H4 & 保边平滑 & \pageref{original:H4}\\
H5 & 八比特与256级 & \pageref{original:H5}\\
H6 & 均衡变换与CDF & \pageref{original:H6}\\
H7 & 同时对比 & \pageref{original:H7}\\
H8 & 混合邻接 & \pageref{original:H8}\\
H9 & 高通核系数和 & \pageref{original:H9}\\
H10 & CLAHE裁剪 & \pageref{original:H10}\\
H11 & 频谱中心化 & \pageref{original:H11}\\
H12 & 理想低通振铃 & \pageref{original:H12}\\
H13 & 周期干扰的陷波 & \pageref{original:H13}\\
\bottomrule
\end{longtable}
